(i) $B_1B_2 = g_2g_3g_3g_1 = g_2g_1 = -B_3$ and $B_2B_1 = B_3$, so $[B_1, B_2] = -2B_3$, and cyclically; $B_k^2 = -1$ and $B_jB_k$ is a bivector for $j \ne k$, so $\langle B_jB_k\rangle_0 = -\delta_{jk}$. A three-dimensional real Lie algebra with these brackets and a negative-definite invariant form is the compact $\su(2)$, the first summand of Theorem \claimref{standard-model.gauge-algebra-minimal-massive-composite}; its simply connected group is the rotor group $\{U : U\widetilde{U} = 1\}$, the unit three-sphere in the coordinates $u_k$. The check verifies the four-component product used for rotors against the Clifford product on all basis pairs.

(ii) The moments of the normalised measure on the unit sphere $S^{n-1}$ are $\prod_k(\alpha_k - 1)!!/\prod_{j < |\alpha|/2}(n + 2j)$ for even $\alpha$ (the Gaussian integral factorised into radial and angular parts), and vanish otherwise by the reflections $u_k \mapsto -u_k$, which preserve the measure; here $n = 4$. \emph{Computed:} the 24 Hurwitz units ($\pm1$, $\pm B_k$, $\tfrac12(\pm1 \pm B_1 \pm B_2 \pm B_3)$) average every monomial of degree at most $5$ to these values, and the eight units $\pm1$, $\pm B_k$ every monomial of degree at most $3$ (not degree $4$: $\tfrac14$ against $\tfrac18$ for $u_0^4$). Every integrand below has degree at most $2$ in each link variable, so either finite group may replace $dU$ link by link.

(iii) $\langle XU\rangle_0$ is linear in $u$, so its average vanishes by (ii). The gluing integrand is bilinear in $u$, and by (ii) $\int u_ku_l\,dU = \tfrac14\delta_{kl}$, so $\int\langle XU\rangle_0\langle\widetilde{U}Y\rangle_0\,dU = \tfrac14\sum_k\langle Xe_k\rangle_0\langle\widetilde{e_k}Y\rangle_0$ over the basis $e_k = 1, B_1, B_2, B_3$; since $\langle\widetilde{e_k}e_l\rangle_0 = \delta_{kl}$, the sum is the expansion of $\langle XY\rangle_0$ in this orthonormal basis. Both sides are bilinear in $(X, Y)$, so the check on all sixteen basis pairs, by the moments and by the Hurwitz average, is complete.

(iv) Along a closed path the inner factors $\widetilde{g_y}g_y = 1$ cancel, leaving $g_xU_C\widetilde{g_x}$, whose scalar part is $\langle U_C\rangle_0$. The Haar measure is invariant under $U \mapsto gU\widetilde{h}$, so for gauge-invariant observables the expectation is unchanged by fixing any tree of links to $1$, and no gauge-fixing term or ghost is needed (the Faddeev-Popov determinant of a tree gauge is $1$); this is Wilson's standard observation for every compact lattice gauge theory, reproduced here for rotor links. The check uses rational rotors at every site of a $2 \times 2$ patch and confirms that an open line is not invariant.

(v) Fix the tree consisting of every horizontal link and the vertical links on the left edge to $1$. Then the plaquette in column $x$, row $y$ is $U_p = V_{x+1,y}\widetilde{V_{x,y}}$ with $V_{0,y} = 1$ and $V_{x,y}$ the vertical link at $x$, and the boundary product is $\langle U_C\rangle_0 = \langle V_{R,0}V_{R,1}\cdots V_{R,T-1}\rangle_0$. In each row the map $(V_{1,y}, \ldots, V_{R,y}) \mapsto (P_{0,y}, \ldots, P_{R-1,y})$, $P_{x,y} = V_{x+1,y}\widetilde{V_{x,y}}$, is a bijection preserving the product Haar measure (left invariance, one factor at a time), and $V_{R,y} = P_{R-1,y}\cdots P_{0,y}$. So the integral is $\int\langle\prod P\rangle_0\prod_{p}\langle P_p\rangle_0$ over $A$ independent Haar rotors, the ordered product running over every plaquette. Integrating one $P_p$ by (iii) with $Y = 1$ (using $\langle P_p\rangle_0 = \langle\widetilde{P_p}\rangle_0$) gives $\tfrac14$ and removes $P_p$ from the product; after $A$ steps the value is $(\tfrac14)^{A}\langle1\rangle_0 = (\tfrac14)^{A}$. In the expansion of the weight, the $A$ distinct plaquettes of the rectangle each taken once come with coefficient $\beta^{A}$, so the leading strong-coupling coefficient is $(\beta/4)^{A}$. \emph{Computed:} $\tfrac14$ for $1 \times 1$ with all four links averaged over the eight units and in the tree gauge over the Hurwitz units, $\tfrac1{16}$ for $2 \times 1$ (Hurwitz) and $\tfrac1{256}$ for $2 \times 2$ (eight units). \emph{Imported} (Wilson 1974; Fact \claimref{synthesis.imported-wilson-lattice-gauge-theory-and-strong-coupling-confinement}(i)): on the infinite lattice no term of lower order contributes, since a plaquette set not covering every link of the loop leaves a link whose average vanishes by (iii), and the normalisation is $1 + O(\beta^2)$; so the string tension at strong coupling is $-\ln(\beta/4) + O(\beta)$. The area law beyond leading order, for all small $\beta$, is the theorem of Osterwalder and Seiler (1978; Fact \claimref{synthesis.imported-wilson-lattice-gauge-theory-and-strong-coupling-confinement}(ii)).

(vi) For $U = a + bB_3$ with $a^2 + b^2 = 1$ and the uniform measure, $\int a^2 = \int b^2 = \tfrac12$ and $\int ab = 0$, so the gluing coefficient on $\mathrm{span}\{1, B_3\}$ is $\tfrac12$, and the argument of (v) gives $(\tfrac12)^{A}$; the four units $\pm1$, $\pm B_3$ reproduce these moments, and the check computes $\tfrac12$ and $\tfrac14$ for the $1 \times 1$ and $2 \times 1$ loops.

(vii) $B^2 = -\|B\|^2$ is a scalar, so $\langle B^{2m}\rangle_0 = (-\|B\|^2)^{m}$ and the odd powers have no scalar part; the check compares the series term by term to sixth order on two rational bivectors.

\emph{Argued, not proved here:} reflection positivity of the Wilson weight (Osterwalder and Seiler; Fact \claimref{synthesis.imported-wilson-lattice-gauge-theory-and-strong-coupling-confinement}(ii)) gives a Hilbert space and a transfer matrix, so (iv) and (v) describe a quantised gauge field at finite spacing; reading $U_p = \exp(B_p)$ with $B_p$ of order the lattice spacing squared times the field strength, (vii) gives the Yang-Mills density with $\beta = 4/g^2$ (Fact \claimref{synthesis.imported-wilson-lattice-gauge-theory-and-strong-coupling-confinement}(i)). By (vi) the area law at strong coupling does not distinguish the non-abelian factor (compact $U(1)$ in four dimensions is proved to deconfine at weak coupling, by Guth (1980) and by Fr\"ohlich and Spencer (1982); Fact \claimref{synthesis.imported-wilson-lattice-gauge-theory-and-strong-coupling-confinement}(iv)), so confinement in the continuum is not decided; the continuum limit at $\beta \to \infty$ needs asymptotic freedom, the sign of the one-loop coefficient, which the result \claimref{synthesis.one-loop-running-coefficients-from-spacetime-algebra-spin-traces}, computes conditionally on the derived content and the doublet count. The same construction with the $\su(3)$ factor uses its own Haar moments ($\tfrac13$ in place of $\tfrac14$ for the normalised trace), not computed here. Transport-absent.
